\begin{tikzpicture}[scale=0.92,>=Stealth,line cap=round,line join=round]
  \path[use as bounding box] (-4.35,-2.05) rectangle (4.35,2.15);

  % 绕过导体的外部场线。
  \draw[blue!70!black,->,thick] (-4.1,1.70) -- (4.1,1.70);
  \draw[blue!70!black,->,thick] (-4.1,-1.70) -- (4.1,-1.70);
  \draw[blue!70!black,->,thick]
    (-4.1,1.03) .. controls (-2.75,1.02) and (-2.10,1.50) ..
    (0,1.56) .. controls (2.10,1.50) and (2.75,1.02) .. (4.1,1.03);
  \draw[blue!70!black,->,thick]
    (-4.1,-1.03) .. controls (-2.75,-1.02) and (-2.10,-1.50) ..
    (0,-1.56) .. controls (2.10,-1.50) and (2.75,-1.02) .. (4.1,-1.03);

  % 一部分场线终止于左侧感应负电荷，另一部分从右侧正电荷出发。
  \foreach \y/\xe in {-0.56/1.43,0/1.65,0.56/1.43}{
    \draw[blue!70!black,->,thick] (-4.1,\y) -- (-\xe,\y);
    \draw[blue!70!black,->,thick] (\xe,\y) -- (4.1,\y);
  }
  \node[blue!70!black,font=\small] at (-3.25,2.02)
    {外电场方向};

  % 导体壳与空腔。
  \fill[gray!22] (0,0) ellipse (1.65 and 1.12);
  \draw[thick] (0,0) ellipse (1.65 and 1.12);
  \fill[white] (0,0) ellipse (0.92 and 0.56);
  \draw[thick] (0,0) ellipse (0.92 and 0.56);

  \foreach \y in {-0.57,-0.19,0.19,0.57}{
    \node[blue!75!black,font=\small] at (-1.48,\y) {$-$};
    \node[red!80!black,font=\small] at (1.48,\y) {$+$};
  }

  \node[font=\small] at (0,1.29) {导体壳};
  \node[font=\small,align=center] at (0,0) {空腔内无电荷\\$E=0$};

  \draw[->,gray!75!black,thick] (-2.20,-1.88) -- (-1.18,-0.82);
  \node[font=\scriptsize,align=center] at (-2.70,-1.90)
    {感应电荷\\重新分布};
  \draw[->,gray!75!black,thick] (2.20,-1.88) -- (1.10,-0.78);
  \node[font=\scriptsize,align=center] at (2.72,-1.90)
    {导体材料内部\\$E=0$};
\end{tikzpicture}
